% !TEX program = xelatex
\documentclass[11pt,a4paper]{ctexart}

\usepackage{amsmath,amssymb,amsfonts}
\usepackage{booktabs}
\usepackage{geometry}
\usepackage{hyperref}
\usepackage{enumitem}
\usepackage{xcolor}
\usepackage{longtable}

\geometry{left=2.5cm,right=2.5cm,top=2.5cm,bottom=2.5cm}
\hypersetup{colorlinks=true,linkcolor=blue,urlcolor=blue}

\title{%
  \textbf{星闪 SLB 物理层}\\[0.4em]
  \Large 6.2.4 物理层控制信息技术文档\\[0.3em]
  \large SAB \textbullet\ RAB \textbullet\ G-PCIB \textbullet\ GCI \textbullet\ TCI \textbullet\ P-PCIB
}
\author{依据 T/XS 10001-2025《星闪无线通信系统 接入层\\
同步低时延宽带空口 SLB 技术要求和测试方法》整理\\
（团体标准 20250326 版）}
\date{2026 年 7 月 3 日}

\begin{document}
\maketitle
\tableofcontents
\newpage

%======================================================================
\part{总览}
%======================================================================

\section{文档范围与模块划分}

本文档对应 SLB 120\,kHz 物理层协议第 6.2.4 节``物理层控制信息''，涵盖 6.2.4.1\textasciitilde 6.2.4.10 全部子节，定义同步/广播/接入信息、G 链路控制信息 GCI、G-PCIB/P-PCIB 控制块及 T 链路控制信息 TCI 的\textbf{比特字段}、\textbf{信道编码与调制}、\textbf{资源映射}与\textbf{盲检规则}。

\begin{table}[h]
\centering
\small
\begin{tabular}{lll}
\toprule
\textbf{子节} & \textbf{内容} & \textbf{核心输出} \\
\midrule
6.2.4.1 & 同步信息 & 通信域/直通链路 80\,bit 同步字段 \\
6.2.4.2 & 同步信息块 SAB & STS/FTS/同步信息 RE 布局 \\
6.2.4.3 & 广播信息 & 72\,bit 物理层配置参数 \\
6.2.4.4 & 接入信息 & 80\,bit 随机接入/调度请求字段 \\
6.2.4.5 & 接入信息块 RAB & T 节点接入块映射与退避 \\
6.2.4.6 & G 链路控制信息 GCI & DSCI 等 7 类 GCI + 编码映射 \\
6.2.4.7 & G-PCIB & CTU 聚合、盲检、表 7 \\
6.2.4.8 & 直通动态调度 & 82\,bit 直通链路 DSCI \\
6.2.4.9 & P-PCIB & 直通控制块符号与映射 \\
6.2.4.10 & T 链路控制信息 TCI & HARQ/CQI/SR 反馈与映射 \\
\bottomrule
\end{tabular}
\end{table}

\section{与 6.2.1/6.2.2/6.2.3/6.2.5 的关系}

6.2.1 将本节 17 类控制信息列入 120\,kHz 系统物理层控制清单，并定义 SAB/RAB/G-PCIB/P-PCIB 等 PIB 打包结构。6.2.2 提供 RE 网格 $(k,l)_p$、超帧/无线帧/TTI 坐标及 G/T/D 符号类型，为本节各控制信息的\textbf{符号索引}（如 \#3、\#5、\#8）提供时频参照。6.2.3 提供 SAB/RAB 内 FTS/STS 序列及 G-DMRS-C、P-DMRS-C、T-DMRS-C 参考信号，与本节控制比特\textbf{同块或同 RE 绑定}映射。6.2.5 物理层数据信息依赖本节 GCI（动态调度数据控制信息、预配置调度激活去激活信息）及 TCI（HARQ/CQI 反馈）完成 TB 调度闭环。

\section{PHY 控制面流水线}

\begin{enumerate}[nosep]
  \item \textbf{下行发现}（6.2.4.1/6.2.4.2/6.2.4.3）：G 节点在 SAB 中映射 STS/FTS/同步信息 $\rightarrow$ 周期广播信息（\#5、\#7）$\rightarrow$ 指示 NSymb-G-PCIB 与 GCI 组位图；
  \item \textbf{接入}（6.2.4.4/6.2.4.5）：T 节点在 RAB 上映射 STS/FTS/接入信息 $\rightarrow$ 可选随路短数据（\#5、\#6）$\rightarrow$ 概率退避重传；
  \item \textbf{G 侧调度}（6.2.4.6/6.2.4.7）：G 节点在 G-PCIB 内按 CTU 聚合映射 GCI $\rightarrow$ T 节点盲检公共/特定 GCI $\rightarrow$ 解析 DSCI/预配置/载波切换等；
  \item \textbf{数据闭环}（6.2.4.6.1 + 6.2.5）：DSCI 指示时频/MCS/HARQ $\rightarrow$ 6.2.5 映射数据 TB $\rightarrow$ T 节点在指示载波反馈 TCI；
  \item \textbf{上行反馈}（6.2.4.10）：TCI（序列或数据方式）映射于 T-DMRS-C 伴随 RE $\rightarrow$ G 节点解调 ACK/CQI/SR；
  \item \textbf{直通链路}（6.2.4.1/6.2.4.8/6.2.4.9）：直通 SAB 指示 NCTU-PCI/NPCI $\rightarrow$ P-PCIB 映射直通 DSCI $\rightarrow$ 6.2.5 直通数据。
\end{enumerate}

\newpage
%======================================================================
\part{6.2.4 物理层控制信息}
%======================================================================

\section{协议章节对应}

本节对应 T/XS 10001-2025 第 6.2.4 节，完整涵盖 6.2.4.1\textasciitilde 6.2.4.10.6。

\section{功能定义}

6.2.4 的核心功能为：在 6.2.2 定义的 RE 网格上承载\textbf{不经过 MAC 复用的物理层控制比特}，支撑：
\begin{enumerate}[nosep]
  \item 网络发现与定时（同步信息、SAB、广播信息）；
  \item 随机接入与资源请求（接入信息、RAB）；
  \item G 节点对 T 节点的调度、预配置、功耗与载波管理（GCI、G-PCIB）；
  \item 定位/感知测量资源配置（G/T 位置测量信号资源指示）；
  \item T 节点闭环反馈（TCI：HARQ、CQI、SR）；
  \item 直通链路动态调度（直通 DSCI、P-PCIB）。
\end{enumerate}

\section{模块边界（上下游及职责划分）}

\subsection{上游模块}

\begin{table}[h]
\centering
\small
\begin{tabular}{p{4cm}p{4.5cm}p{5.5cm}}
\toprule
\textbf{上游} & \textbf{输入} & \textbf{说明} \\
\midrule
6.2.2 帧结构 & 符号类型、$(k,l)$ 网格 & RE 坐标与 TTI 长度 \\
6.2.3 物理层信号 & FTS/STS/DMRS 序列 & SAB/RAB/TCI 同块信号 \\
高层 RRC/XRC & ACK 资源池、DRX、预配置 ID 等 & GCI/TCI 高层语义 \\
6.2.6 信道编码 & $C=1$ 极化码参数 & 控制信息编码 \\
6.5.1/6.5.6 & CRC 多项式、QPSK 调制 & 比特$\rightarrow$符号 \\
\bottomrule
\end{tabular}
\end{table}

\subsection{下游模块}

\begin{table}[h]
\centering
\small
\begin{tabular}{p{4cm}p{4.5cm}p{5.5cm}}
\toprule
\textbf{下游} & \textbf{本模块输出} & \textbf{说明} \\
\midrule
6.2.5 数据信息 & DSCI/预配置指示的资源与 MCS & TB 调度 \\
接收机盲检 & G-PCIB 上 CTU 候选 & GCI 解析 \\
MAC HARQ & TCI ACK/NACK、CQI & 重传与链路自适应 \\
7.2 媒体接入层 & 随机接入、SR 流程 & 7.2.2.1/7.2.2.2 \\
6.6/6.7 测量 & 位置/感知资源指示 & PMS 时频配置 \\
\bottomrule
\end{tabular}
\end{table}

\subsection{本模块不负责的内容}

物理层数据 TB 比特语义与 MCS 表（6.2.5）、信号序列生成（6.2.3）、极化码内核（6.2.6 算法细节）、OFDM 合成（6.2.9）、射频指标（6.2.10/第 8 章）。

\section{在 SLB 通信流程中的角色}

6.2.4 处于 PHY \textbf{控制面承载层}：SAB 完成粗同步与系统参数广播 $\rightarrow$ RAB/G-PCIB 完成接入与调度 $\rightarrow$ 6.2.5 在 GCI 指示资源上映射数据 $\rightarrow$ TCI 闭合 HARQ/CQI 环路。直通场景下 P-PCIB 替代 G-PCIB 承担同类调度语义。

%----------------------------------------------------------------------
\section{强制要求与关键参数（T/XS 10001-2025 原文）}
%----------------------------------------------------------------------

\subsection{6.2.4.1 同步信息}

\subsubsection{6.2.4.1.1 信息比特}

同步信息共 80 比特。

从最低位比特到最高位比特，通信域同步信息中，有效比特 52 比特，预留比特 4 比特，CRC24 比特，具体包含如下信息：
\begin{itemize}[nosep]
  \item 3 比特：信息格式指示，对于通信域同步信息，该字段取值为 0。
  \item 24 比特：G 节点标识。G 节点开机后随机产生。
  \item 8 比特：G 节点多域同步能力属性。定义见 6.2.10 节。
  \item 2 比特：同步信息块传输周期，0：无周期，1：2 超帧，2：4 超帧，3：8 超帧。同步信息块传输周期为无周期，仅用于非连续传输模式，此时信道占用期间仅发送 1 次同步信息块，信道占用时长为整数个无线帧。
  \item 4 比特：（包括当前同步信息块周期）剩余同步信息块周期数量。等于 15 时为连续传输模式，等于 14 时为连续传输模式结尾带 RACH，等于 0 时为非连续传输模式结尾带 RACH。
  \item 8 比特：指示本基础载波包含哪些 T 节点组的 GCI。使用位图方式指示，其中位图的最低位比特指示本基础载波是否包含公共 GCI，其它每一比特对应 1 个 T 节点组。对于本字段的每 1 个比特，0：不包含对应 T 节点组的 GCI，1：包含对应 T 节点组的 GCI。
  \item 1 比特：符号类型指示。0：类型一 A 或类型二；1：类型三、类型四或类型一 B。
  \item 1 比特：发送同步信息的基础信道是否允许直通链路使用。0：不允许，1：允许。
  \item 1 比特：广播信息存在性指示。0：发送当前同步信息的 TTI 中不发送广播信息；1：发送当前同步信息的 TTI 中发送广播信息。
  \item 4 比特：预留，当前版本全部为 0。
  \item 24 比特：根据本文件 6.5.1 节的处理方法，使用循环冗余校验生成多项式 $g_{\mathrm{CRC24B}}(D)$ 计算循环冗余校验序列。
\end{itemize}

从最低位比特到最高位比特，直通链路同步信息中，有效比特 55 比特，预留比特 4 比特，CRC24 比特，具体包含以下信息：
\begin{itemize}[nosep]
  \item 3 比特：信息格式指示。对于直通链路同步信息，本字段取值为 1。
  \item 24 比特：先发节点标识。
  \item 24 比特：后发节点标识，或后发节点组标识。当前 COT 专用于服务发现时该字段取值为 1。
  \item 2 比特：直通链路控制信息聚合级别 NCTU-PCI。0：聚合级别 NCTU-PCI=1，1：聚合级别 NCTU-PCI=2，2：聚合级别 NCTU-PCI=4，3：聚合级别 NCTU-PCI=8。直通链路控制信息聚合级别的定义与 G 链路控制信息聚合级别相同，参考 6.2.4.7 节。
  \item 2 比特：直通链路控制信息数量 NPCI。数值等于直通链路控制信息的数量。直通链路控制信息的数量大于 1 时，所有直通链路控制信息的聚合级别相同，都与直通链路控制信息聚合级别字段的指示一致。
  \item 1 比特：预留，当前版本为 0。
  \item 24 比特：根据本文件 6.5.1 节的处理方法，使用循环冗余校验生成多项式 $g_{\mathrm{CRC24B}}(D)$ 计算循环冗余校验序列。
\end{itemize}

\subsubsection{6.2.4.1.2 信道编码和调制}

同步信息使用本文件 6.2.6 规定的信道编码参数，进行 6.2.6.1.1、6.2.6.1.2、6.2.6.1.3 和 6.2.6.1.4 的极化编码，产生比特序列 $g_k$，其中 $k = 0, \ldots, G-1$，并根据本文件 6.5.6 节的方法使用 QPSK 调制将比特序列 $g_k$ 映射为调制符号序列 $d_i$。

\subsubsection{6.2.4.1.3 资源映射}

当无线帧中传输同步信息时，同步信息编码调制后的符号序列 $d_i$ 在该无线帧的 \#3 符号和 \#4 符号共两个符号中按照符号的时间顺序，在每个符号中按照先低频率子载波再高频率子载波的顺序进行映射。

\subsection{6.2.4.2 同步信息块 SAB}

在连续传输模式下，从 \#0 无线帧开始，以 [2 毫秒，4 毫秒，8 毫秒] 为周期出现同步信息块。

在非连续传输模式下，在占用的部分或全部基础载波上，从 COT 第一个无线帧开始，以 [2 毫秒，4 毫秒，8 毫秒] 为周期出现同步信息块，一直到该组信道占用无线帧结束。

在有同步信息块的基础载波上，有同步信息块的无线帧中，\#0，\#1，\#2，\#3，\#4 符号为同步信息块的资源。

对于通信域上的同步信息块，通信域 G 节点在同步信息块中 \#0 符号映射 STS，\#1，\#2 符号映射 FTS，\#3，\#4 符号映射通信域同步信息。其中，STS 序列，$u = 1, 2, \ldots, 16$，为通信域第二训练信号标识，其值为通信域 G 节点标识最低 4 位对应的数值加 1；FTS 序列，$u = 1, 160$ 为通信域同步信息块第一训练信号标识，与符号类型对应，$u = 1$ 对应类型一 A、类型一 B 或类型四，$u = 160$ 对应类型二或类型三。

对于直通链路上的同步信息块，直通链路先发节点在同步信息块中 \#0 符号映射 STS，\#1，\#2 符号映射 FTS，\#3，\#4 符号映射直通链路同步信息。其中，STS 序列 $u = 1, 2, \ldots, 16$，为直通链路第二训练信号标识，其值为直通链路先发节点标识最低 4 位对应的数值加 1；FTS 序列 $u = 159$ 为直通链路同步信息块第一训练信号标识。

在通信域或直通链路上，若存在无同步信息块 SAB 的基础载波，那么在至少一个基础载波上有同步信息块的无线帧中，无同步信息块 SAB 的基础载波上该无线帧 \#0 符号映射 STS 序列。该 STS 序列的 $u = 1, 2, \ldots, 16$，为通信域第二训练信号标识，其值为通信域 G 节点标识最低 4 位对应的数值加 1。

\subsection{6.2.4.3 广播信息}

\subsubsection{6.2.4.3.1 信息比特}

广播信息中，有效比特 37 比特，预留比特 11 比特，CRC24 比特，共 72 比特，承载物理层配置参数，从最低位比特到最高位比特，具体包含如下信息：
\begin{itemize}[nosep]
  \item 3 比特：符号类型指示。0：类型一 A，1：类型一 B，2：类型二，3：类型三，4：类型四，其它取值预留。
  \item 1 比特：传输模式信息。0：连续传输模式，1：非连续传输模式。
  \item 19 比特：无线帧索引号。指示发送该广播信息的第一个符号所在无线帧的索引号。无线帧索引号的高 16 位为无线帧所在超帧的超帧编号，无线帧索引号的低 3 位为该无线帧的无线帧编号。
  \item 3 比特：TTI 长度指示，0：1 无线帧，1：2 无线帧，2：4 无线帧，3：8 无线帧，4：16 无线帧，5：32 无线帧，6：64 无线帧，7：预留。
  \item 3 比特：通信域第一组系统消息 SIB0 的发送周期，0 指示可能发送系统消息的帧的起始编号为 64 的倍数，1 指示可能发送系统消息的帧的起始编号为 128 的倍数，2 指示可能发送系统消息的帧的起始编号为 256 的倍数，3 指示可能发送系统消息的帧的起始编号为 512 的倍数；4 指示可能发送系统消息的帧的起始编号为 1024 的倍数，5 指示可能发送系统消息的帧的起始编号为 2048 的倍数，6 指示可能发送系统消息的帧的起始编号为 4096 的倍数，7 指示可能发送系统消息的帧的起始编号为 8192 的倍数。
  \item 2 比特：指示通信域第一组系统消息 SIB0 在一个发送周期内可能使用的连续的 TTI 长度的个数，该字段指示的连续的 TTI 中的公共 GCI 资源中发送 SIB0 的调度信息，0 指示从起始超帧编号开始最多在 8 个连续的 TTI 长度内发送系统消息，1 指示从起始超帧编号开始最多在 16 个连续的 TTI 长度内发送系统消息，2 指示从起始超帧编号开始最多在 32 个连续的 TTI 长度内发送系统消息，3 指示从起始超帧编号开始最多在 64 个连续的 TTI 长度内发送系统消息。
  \item 3 比特：G 链路物理层控制信息块 G-PCIB 符号个数信息 NSymb-G-PCIB，0：NSymb-G-PCIB=1，1：NSymb-G-PCIB=2，2：NSymb-G-PCIB=3，3：NSymb-G-PCIB=4，4：NSymb-G-PCIB=6，5：NSymb-G-PCIB=8，6：NSymb-G-PCIB=10，7：NSymb-G-PCIB=12。在包含同步信息块 SAB 和物理层广播信息的 TTI 中，G-PCIB 为从 \#8 符号开始的连续 NSymb-G-PCIB 个符号；在包含同步信息块 SAB 且不包含物理层广播信息的 TTI 中，若 NSymb-G-PCIB$>1$ 则 G-PCIB 为 \#5 符号和从 \#7 符号开始的连续 $N-1$ 个符号，若 NSymb-G-PCIB=1 则 G-PCIB 仅包含 \#5 符号；在其它 TTI 中，G-PCIB 为从 \#1 符号开始的连续 NSymb-G-PCIB 个符号。
  \item 1 比特：广播信息或系统消息变更指示信息，0 指示广播信息和系统消息在下一个变更周期内不发生变更，1 指示广播信息或系统消息在下一个变更周期内发生变更。
  \item 2 比特：广播信息或系统消息的变更周期，0 指示 512 无线帧，1 指示 1024 无线帧，2 指示 2048 无线帧，3 指示 4096 无线帧。
  \item 11 比特：预留比特，当前版本全部为 0。
  \item 24 比特：根据本文件 6.5.1 节的处理方法，使用循环冗余校验生成多项式 $g_{\mathrm{CRC24B}}(D)$ 计算循环冗余校验序列。
\end{itemize}

\subsubsection{6.2.4.3.2 信道编码和调制}

广播信息使用本文件 6.2.6 规定的信道编码参数，进行 6.2.6.1.1、6.2.6.1.2、6.2.6.1.3 和 6.2.6.1.4 的极化编码，产生比特序列 $g_k$，其中 $k = 0, \ldots, G-1$，并根据本文件 6.5.6 节的方法使用 QPSK 调制将比特序列 $g_k$ 映射为调制符号序列 $d_i$。

\subsubsection{6.2.4.3.3 资源映射}

在连续传输模式下，G 节点从 \#0 无线帧开始，以 8 毫秒为周期发送广播信息。

在非连续传输模式下，在每一个 COT 中，从 COT 的第一个无线帧开始，以 8 毫秒为周期发送广播信息，一直到该 COT 结束。

当无线帧中传输广播信息时，广播信息编码调制后的符号序列 $d_i$ 在该无线帧的 \#5 符号和 \#7 符号共两个符号中按照符号的时间顺序，在每个符号中按照先低频率子载波再高频率子载波的顺序进行映射。

\subsection{6.2.4.4 接入信息}

\subsubsection{6.2.4.4.1 信息比特}

接入信息中，有效比特 56 比特，CRC24 比特，共 80 比特，从最低位比特到最高位比特，具体包含如下信息：
\begin{itemize}[nosep]
  \item 3 比特：信息格式指示，对于接入信息，该字段取值为 2。
  \item 24 比特：T 节点标识，T 节点随机产生。
  \item 24 比特：G 节点标识，与待接入 G 节点发送的同步信息中的 G 节点标识一致。
  \item 3 比特：接入目的。0：初始接入，1：调度请求，2：T 节点请求 G 节点配置以进行无连接测量，但不接入 G 节点，3：服务发现，4：随路短数据传递。
  \item 2 比特：指示初始接入阶段 T 节点带宽类型。0：初始接入阶段 T 节点工作载波包含 1 个基础载波，1：初始接入阶段 T 节点工作载波包含 4 个基础载波，2：初始接入阶段 T 节点工作载波包含本频段内所有基础载波，3：预留。
  \item 24 比特：根据本文件 6.5.1 节的处理方法，使用循环冗余校验生成多项式 $g_{\mathrm{CRC24B}}(D)$ 计算循环冗余校验序列。
\end{itemize}

\subsubsection{6.2.4.4.2 信道编码和调制}

接入信息使用本文件 6.2.6 规定的信道编码参数，进行 6.2.6.1.1、6.2.6.1.2、6.2.6.1.3 和 6.2.6.1.4 的极化编码，产生比特序列 $g_k$，其中 $k = 0, \ldots, G-1$，并根据本文件 6.5.6 节的方法使用 QPSK 调制将比特序列 $g_k$ 映射为调制符号序列 $d_i$。

\subsubsection{6.2.4.4.3 资源映射}

当无线帧中传输接入信息时，接入信息编码调制后的符号序列 $d_i$ 在该无线帧的 \#3 符号和 \#4 符号共两个符号中按照符号的时间顺序，在每个符号中按照先低频率子载波再高频率子载波的顺序进行映射。

\subsection{6.2.4.5 接入信息块 RAB}

在连续传输模式下，接入信息块的资源由系统消息配置，可配置为满足条件的同步信息块发送周期的最后一个无线帧。

在非连续传输模式下，在有同步信息块的基础载波上，一个 COT 最后 1 个无线帧可以配置为接入信息块的资源。COT 最后 1 个无线帧中是否配置接入信息块资源，在该 COT 最后一个同步信息块的同步信息中指示。

在有接入信息块的基础载波上，有接入信息块的无线帧中，\#0，\#1，\#2，\#3，\#4 符号为 T 节点映射接入信息块对应的物理层控制信号和物理层控制信息的资源。

T 节点在其中 \#0 符号映射 STS，\#1，\#2 符号映射 FTS，\#3，\#4 符号映射接入信息。其中，STS 序列，$u = 1, 2, \ldots, 16$，为通信域第二训练信号标识，其值为通信域 G 节点标识最低 4 位对应的数值加 1；FTS 序列，$u = 2$ 为通信域接入信息块第一训练信号标识。

当接入信息中``接入目的''字段取值为``4''时，接入信息块还包括本基础载波上包含本接入信息块的无线帧中的 \#5，\#6 符号。其中 \#5，\#6 符号映射高层数据。高层数据包括 48 比特有效数据和 24 比特 CRC，使用本文件 6.2.6 规定的信道编码参数，进行 6.2.6.1.1、6.2.6.1.2、6.2.6.1.3 和 6.2.6.1.4 的极化编码，产生比特序列 $g_k$，其中 $k = 0, \ldots, G-1$，并根据本文件 6.5.6 节的方法使用 QPSK 调制将比特序列 $g_k$ 映射为调制符号序列 $d_i$。

T 节点基于上行定时映射接入信息块对应的物理层控制信号和物理层控制信息，若 T 节点无法确定上行定时，则以下行定时作为上行定时映射接入信息块对应的物理层控制信号和物理层控制信息。

T 节点以概率 $P$ 在接入信息块映射对应的物理层控制信号和物理层控制信息。$P$ 的初始值为 1；每当 T 节点在接入信息块上映射对应的物理层控制信号和物理层控制信息后，在下一次检测到同一个 G 节点的 COT 中未收到对应的响应信息，$P$ 更新为当前值的一半；每当 T 节点在接入信息块上映射对应的物理层控制信号和物理层控制信息后，在下一次检测到同一个 G 节点的 COT 中收到对应的响应信息，$P$ 更新为 1。

\subsection{6.2.4.6 G 链路控制信息 GCI}

G 链路控制信息，包括动态调度数据控制信息、预配置调度激活去激活信息、休眠与唤醒指示信息、快速载波切换指示信息、初始接入载波指示信息、G 链路位置测量信号资源指示信息、T 链路位置测量信号资源指示信息。

G 链路控制信息 GCI 分为公共 GCI 和 T 节点特定 GCI 两类，其中 T 节点特定 GCI 的循环冗余校验序列使用 T 节点物理层标识加扰，公共 GCI 的循环冗余校验序列使用其它方式加扰。

\subsubsection{6.2.4.6.1 动态调度数据控制信息}

G 节点通过高层信令配置 1 个 TB 包含的最大 CBG 个数为 $N$。在一次传输中，一个 TB 包含的 CB 为 $C$，则该 TB 实际包含的 CBG 的数为 $M = \min(C, N)$。$M_1 = \mod(C, M)$，如果 $M_1 > 0$，$K_1 = \lceil C/M \rceil$，$K_2 = \lfloor C/M \rfloor$，在 $M$ 个 CBG 中，\#0 到 \#$(M_1-1)$ 的 CBG 包含 $K_1$ 个 CB，\#$M_1$ 到 \#$(M-1)$ 的 CBG 包含 $K_2$ 个 CB；如果 $M_1 = 0$，每个 CBG 包含 $(C/M)$ 个 CB。

当高层 \texttt{schedMode} 和 \texttt{harqFeedBackMode} 配置被调度的 T 节点基于传输块 TB 的重传并使用基于 TB 的 HARQ 反馈时，或高层未配置 \texttt{schedMode} 和 \texttt{harqFeedBackMode}，所以默认基于传输块 TB 的重传和基于 TB 的 HARQ 反馈时，又或者调度广播数据时，动态调度数据控制信息有效比特 58 比特，CRC24 比特，共 82 比特，从最低位比特到最高位比特，包括：
\begin{itemize}[nosep]
  \item 4 比特：物理层控制信息格式指示。对于基于传输块 TB 的重传并进行基于 TB 的 HARQ 反馈，该字段的取值为 0。
  \item 1 比特：链路类型指示信息。0：G 链路传输，1：T 链路传输。
  \item 3 比特：解调参考信号端口数量指示信息。解调参考信号端口数量是（3 比特数值+1）。
  \item 16 比特：频域资源指示信息。16 比特中从最低位的比特开始到最高位的比特结束，与节点工作子载波组从小到大的顺序一一对应。16 比特中，取值为 1 的比特代表使用该比特对应的节点工作子载波组，取值为 0 的比特代表不使用该比特对应的节点工作子载波组。初始接入阶段 XRC 重配置成功之前，本字段中节点工作子载波组为 T 节点工作子载波组，XRC 重配置成功之后，本字段中节点工作子载波组为 G 节点工作子载波组。当本信息用于调度广播数据，且本字段的指示范围为发送本信息的基础载波时，本字段中节点工作子载波组为发送本信息的基础载波上的基础子载波组。
  \item 7 比特：起始符号指示。在该 TTI 内，起始符号的索引。对于小于或等于 1\,ms 的 TTI，其指示单位为 1 个符号；对于 2\,ms 的 TTI，其指示单位为 2 个符号；对于 4\,ms TTI，其指示单位为 4 个符号；对于 8\,ms TTI，其指示单位为 8 个符号。
  \item 7 比特：结束符号指示。在该 TTI 内，结束符号的索引。对于小于或等于 1\,ms 的 TTI，其指示单位为 1 个符号；对于 2\,ms 的 TTI，其指示单位为 2 个符号；对于 4\,ms TTI，其指示单位为 4 个符号；对于 8\,ms TTI，其指示单位为 8 个符号。
  \item 2 比特：HARQ 进程指示。
  \item 1 比特：新包指示。当该值翻转时，指示传输新包，当该值不变时，指示重传旧包。
  \item 5 比特：调制编码方式指示信息。
  \item 3 比特：调制符号重复指示。对于单流传输，该字段指示该单流传输的重复次数；对于两流传输，该字段指示第 2 流传输时符号的重复次数。该字段通过指示重复次数集合中元素的索引来确定重复次数，重复次数集合根据 XRC 重配置信令配置。调度系统消息时以及未配置时，重复次数由该字段数值+1 确定。
  \item 2 比特：冗余版本指示。0 指示冗余版本 0，1 指示冗余版本 1，2 指示冗余版本 2，3 指示冗余版本 3。
  \item 3 比特：在高层 XRC 信令配置的 8 个 ACK/NACK 资源中指示当前使用的资源。当指示的资源为空时，当前进程的 ACK/NACK 缓存并不反馈；当指示的资源非空时，按照进程号升序反馈当前及缓存的进程的 ACK/NACK。当动态调度数据控制信息用于调度广播数据（例如 G 链路系统消息）时，以及用于调度 T 链路数据时，本字段无效。当动态调度数据控制信息用于调度 XRC 建立消息 XRCSetup 时，本字段无效。
  \item 4 比特：TCI 所在的工作信道内的基础载波序号指示。单级 HARQ 反馈信息或两级 HARQ 反馈信息所在的工作信道内的基础载波序号（4 比特数值+1）。当动态调度数据控制信息用于调度广播数据（例如 G 链路系统消息）时，以及用于调度 T 链路数据时，本字段无效。当动态调度数据控制信息用于调度 XRC 建立消息 XRCSetup 时，本字段无效。
  \item 24 比特：根据本文件 6.5.1 节的处理方法，使用循环冗余校验生成多项式 $g_{\mathrm{CRC24B}}(D)$ 计算循环冗余校验序列，并使用 24 比特物理层标识进行加扰。当动态调度数据控制信息用于调度广播数据（例如 G 链路系统消息）时，该 24 比特物理层标识为 G 节点物理层标识。当本信息用于调度广播数据，且``频域资源指示信息''字段的指示范围为发送本信息的基础载波时，该 24 比特物理层标识为``0''。
\end{itemize}

当高层 \texttt{schedMode} 和 \texttt{harqFeedBackMode} 配置被调度的 T 节点支持基于编码块组 CBG 的重传并进行基于 CBG 的 HARQ 反馈时，动态调度数据控制信息有效比特 62 比特，预留比特 2 比特，CRC24 比特，动态调度数据控制信息共 88 比特，从最低位比特到最高位比特，包括：
\begin{itemize}[nosep]
  \item 1 比特：链路类型指示信息。0：G 链路传输，1：T 链路传输。
  \item 3 比特：解调参考信号端口数量指示信息。解调参考信号端口数量是（3 比特数值+1）。
  \item 16 比特：频域资源指示信息。16 比特中从最低位的比特开始到最高位的比特结束，与 G 节点工作子载波组从小到大的顺序一一对应。16 比特中，取值为 1 的比特代表使用该比特对应的 G 节点工作子载波组，取值为 0 的比特代表不使用该比特对应的 G 节点工作子载波组。
  \item 7 比特：起始符号指示。在该 TTI 内，起始符号的索引。对于小于或等于 1\,ms 的 TTI，其指示单位为 1 个符号，对于 2\,ms 的 TTI，其指示单位为 2 个符号，对于 4\,ms TTI，其指示单位为 4 个符号；对于 8\,ms TTI，其指示单位为 8 个符号。
  \item 7 比特：结束符号指示。在该 TTI 内，结束符号的索引。对于小于或等于 1\,ms 的 TTI，其指示单位为 1 个符号，对于 2\,ms 的 TTI，其指示单位为 2 个符号，对于 4\,ms TTI，其指示单位为 4 个符号；对于 8\,ms TTI，其指示单位为 8 个符号。
  \item 2 比特：HARQ 进程指示。
  \item 1 比特：新包指示。当该值翻转时，指示传输新包，当该值不变时，指示重传旧包。
  \item 5 比特：调制编码方式指示信息。
  \item 3 比特：调制符号重复指示。对于单流传输，该字段指示该单流的重复次数；对于两流传输，该字段指示第 2 流传输时符号的重复次数。该字段通过指示重复次数集合中元素的索引来确定重复次数，重复次数集合根据 XRC 重配置信令配置。调度系统消息时以及未配置时，重复次数由该字段数值+1 确定。
  \item 2 比特：冗余版本指示。0：冗余版本 0，1：冗余版本 1，2：冗余版本 2，3：冗余版本 3。
  \item 3 比特：在高层 XRC 信令配置的 8 个 ACK/NACK 资源中指示当前使用的资源。当指示的资源为空时，当前进程的 ACK/NACK 缓存并不反馈；当指示的资源非空时，按照进程号升序反馈当前及缓存的进程的 ACK/NACK。当动态调度数据控制信息用于调度 T 链路数据时，本字段无效。
  \item 4 比特：TCI 所在的工作信道内的基础载波序号指示。单级 HARQ 反馈信息或两级 HARQ 反馈信息所在的工作信道内的基础载波序号（4 比特数值+1）。当动态调度数据控制信息用于调度 T 链路数据时，本字段无效。
  \item $X$ 比特：以位图的形式指示重传的 CBG。$X$ 由高层 \texttt{singleStageCBGNumber} 配置，表示 TB 块最大可包含的 CBG 个数，$X$ 可以为 2，4，8。调度重传时，$X$ 比特中值为 0 的比特代表该比特对应的 CBG 不进行重传，值为 1 的比特代表该比特对应的 CBG 进行重传。调度初传时，每个比特的值都为 1。
  \item $Y$ 比特：填充比特，全部为 0。$X+Y=8$。
  \item 2 比特：预留比特，当前版本全部为 0。
  \item 24 比特：根据本文件 6.5.1 节的处理方法，使用循环冗余校验生成多项式 $g_{\mathrm{CRC24B}}(D)$ 计算循环冗余校验序列，并使用 24 比特物理层标识进行加扰。
\end{itemize}

当高层 \texttt{harqFeedBackMode} 配置被调度的 T 节点对 G 链路数据采用两级 HARQ 反馈时，或当高层 \texttt{schedMode} 配置被调度的 T 节点采用两级动态调度数据控制信息调度 T 链路数据时，第一级动态调度数据控制信息有效比特 65 比特，预留比特 2 比特，CRC 24 比特，动态调度数据控制信息共 91 比特，从最低位比特到最高位比特，包括：
\begin{itemize}[nosep]
  \item 1 比特：GCI 格式指示。本信息本字段为 0。
  \item 1 比特：链路类型指示信息。0：G 链路传输，1：T 链路传输。
  \item 3 比特：解调参考信号端口数量指示信息。解调参考信号端口数量是（3 比特数值+1）。
  \item 16 比特：频域资源指示信息。16 比特中从最低位的比特开始到最高位的比特结束，与 G 节点工作子载波组从小到大的顺序一一对应。16 比特中，取值为 1 的比特代表使用该比特对应的 G 节点工作子载波组，取值为 0 的比特代表不使用该比特对应的 G 节点工作子载波组。
  \item 7 比特：起始符号指示。在该 TTI 内，起始符号的索引。对于小于或等于 1\,ms 的 TTI，其指示单位为 1 个符号，对于 2\,ms 的 TTI，其指示单位为 2 个符号，对于 4\,ms TTI，其指示单位为 4 个符号；对于 8\,ms TTI，其指示单位为 8 个符号。
  \item 7 比特：结束符号指示。在该 TTI 内，结束符号的索引。对于小于或等于 1\,ms 的 TTI，其指示单位为 1 个符号，对于 2\,ms 的 TTI，其指示单位为 2 个符号，对于 4\,ms TTI，其指示单位为 4 个符号；对于 8\,ms TTI，其指示单位为 8 个符号。
  \item 2 比特：HARQ 进程指示。
  \item 1 比特：新包指示。当该值翻转时，指示传输新包，当该值不变时，指示重传旧包。
  \item 5 比特：调制编码方式指示信息。
  \item 3 比特：调制符号重复指示。对于单流传输，该字段指示该单流的重复次数；对于两流传输，该字段指示第二流的重复次数。该字段通过指示重复次数集合中元素的索引来确定重复次数，重复次数集合根据 XRC 重配置信令配置。调度系统消息时以及未配置时，重复次数由该字段数值+1 确定。
  \item 2 比特：冗余版本指示。0：冗余版本 0，1：冗余版本 1，2：冗余版本 2，3：冗余版本 3。
  \item 3 比特：在高层 XRC 信令配置的 8 个 ACK/NACK 资源中指示当前使用的资源。当指示的资源为空时，当前进程的 ACK/NACK 缓存并不反馈；当指示的资源非空时，按照进程号升序反馈当前及缓存的进程的 ACK/NACK。当动态调度数据控制信息用于调度 T 链路数据时，本字段无效。
  \item 4 比特：TCI 所在的工作信道内的基础载波序号指示。单级 HARQ 反馈信息或两级 HARQ 反馈信息所在的工作信道内的基础载波序号（4 比特数值+1）。当动态调度数据控制信息用于调度 T 链路数据时，本字段无效。
  \item $X$ 比特：以位图的形式指示重传的 CBG。$X$ 由高层 \texttt{singleStageCBGNumber} 配置，表示 TB 块最大可包含的 CBG 个数，$X$ 可以为 2，4，8。调度重传时，$X$ 比特中值为 0 的比特代表该比特对应的 CBG 不进行重传，值为 1 的比特代表该比特对应的 CBG 进行重传。调度初传时，每个比特的值都为 1。
  \item $Y$ 比特：填充比特，全部为 0。$X+Y=8$。
  \item 2 比特：重传类型指示。0：通过第一级动态调度数据控制信息指示的 CBG 重传，1：通过第一级和第二级动态调度数据控制信息指示 T 链路的 CBG 重传，2：基于两级 HARQ 反馈信息中的第一级反馈的 G 链路的重传，3：基于两级 HARQ 反馈信息中的第一级和第二级反馈的 G 链路的重传。当重传类型指示的值为 0 或者 1 时，上述 $X$ 比特和 $Y$ 比特有效，否则 $X$ 比特和 $Y$ 比特的值都为 0。调度初传时，本字段无效。
  \item 2 比特：预留比特，当前版本全部为 0。
  \item 24 比特：根据本文件 6.5.1 节的处理方法，使用循环冗余校验生成多项式 $g_{\mathrm{CRC24B}}(D)$ 计算循环冗余校验序列，并使用 24 比特物理层标识进行加扰。
\end{itemize}

当高层 \texttt{schedMode} 配置被调度的 T 节点采用两级动态调度数据控制信息调度 T 链路数据，且对应第一级动态调度数据控制信息中``重传类型指示''字段为``1''时，第二级动态调度数据控制信息有效比特 67 比特，CRC24 比特，动态调度数据控制信息共 91 比特，从最低位比特到最高位比特，包括：
\begin{itemize}[nosep]
  \item 1 比特：GCI 格式指示。本信息本字段为 1。
  \item 64 比特：以位图的形式指示重传的子 CBG。
  \item 2 比特：预留比特，当前版本全部为 0。
  \item 24 比特：根据本文件 6.5.1 节的处理方法，使用循环冗余校验生成多项式 $g_{\mathrm{CRC24B}}(D)$ 计算循环冗余校验序列，并使用 24 比特物理层标识进行加扰。
\end{itemize}

第二级动态调度数据控制信息不用于 G 链路数据的调度，也不用于初传调度。

第二级动态调度数据控制信息包含的子 CBG 的个数 $N_2$ 为 64。第二级动态调度数据控制信息指示中，每个子 CBG 对应的 CB 根据第一级动态调度数据控制信息指示的错误的 CBG 包含的 CB 到子 CBG 的映射确定。根据第一级动态调度数据控制信息，一个 TB 中错误的 CBG 包含的 CB 个数为 $C$，则该错误的 CBG 实际包含的子 CBG 的个数为 $M = \min(C, N_2)$。$M_1 = \mod(C, M)$，如果 $M_1 > 0$，$K_1 = \lceil C/M \rceil$，$K_2 = \lfloor C/M \rfloor$，在 $M$ 个子 CBG 中，\#0 到 \#$(M_1-1)$ 的子 CBG 包含 $K_1$ 个 CB，\#$M_1$ 到 \#$(M-1)$ 的子 CBG 包含 $K_2$ 个 CB；如果 $M_1 = 0$，每个子 CBG 包含 $(C/M)$ 个 CB。

T 节点在接入 G 节点过程中使用基于传输块 TB 的调度并进行基于 TB 的 HARQ 反馈。G 节点可以通过 XRC 重配置过程配置 T 节点使用基于 CBG 的重传并进行基于 CBG 的 HARQ 反馈，或者配置 T 节点对 G 链路数据采用两级 HARQ 反馈并采用两级动态调度数据控制信息调度 T 链路数据。G 节点对 T 节点进行重配置时，需重新配置 24 比特物理层标识 Phy-ID，不同 Phy-ID 对应不同 T 节点的行为。G 节点接收到 T 节点重配置完成信息后，采用新的 Phy-ID 进行调度，T 节点对于新旧两个不同 Phy-ID 的处理取决于 T 节点的实现。

\subsubsection{6.2.4.6.2 预配置调度激活去激活信息}

G 节点可以通过高层信令给 T 节点分配多个预配置调度标识以及每个标识对应的预配置调度的配置信息。T 节点接收并存储配置信息以及相应的标识，但不激活对应的传输。G 节点通过发送预配置调度激活或者去激活 T 节点存储的预配置调度，并按照对应的配置传输物理层数据信息、物理层控制信息、物理层信号或测量。

预配置调度适用于 G 链路、T 链路、附链路，可用于单播、组播和广播，可以配置单次传输、周期传输。对于非连续传输模式，COT 结束则当前激活的预配置调度去激活。

预配置调度激活去激活信息中，有效比特 52 比特，预留比特 6 比特，CRC24 比特，共 82 比特。从最低位比特到最高位比特，预配置调度激活去激活信息包括：
\begin{itemize}[nosep]
  \item 16 比特：频域资源指示。16 比特中从最低位的比特开始到最高位的比特结束，与 G 节点工作子载波组从小到大的顺序一一对应。16 比特中，取值为 1 的比特代表使用该比特对应的 G 节点工作子载波组，取值为 0 的比特代表不使用该比特对应的 G 节点工作子载波组。指示激活的预配置调度的频域资源。若本信息调度数据，则本字段指示的使用工作子载波组的数量与对应高层信令指示的传输数据使用的工作子载波组的数量一致。若本信息调度信号，则本字段指示信号传输使用的工作子载波组，在信号传输使用的工作子载波组内，信号传输是否使用梳齿、使用哪个或哪些梳齿、梳齿与天线端口对应方式等，由高层配置。若本信息仅用于去激活，则该字段为 0。
  \item 16 比特：激活预配置调度的预配置调度标识。若本信息仅用于去激活，则该字段指示未配置的预配置调度标识。
  \item 16 比特：去激活预配置调度的预配置调度标识。若本信息仅用于激活，则该字段指示未配置的预配置调度标识。
  \item 4 比特：TCI 所在的工作信道内的基础载波序号指示。本信息对应的 TCI 反馈所在的工作信道内的基础载波序号（4 比特数值+1）。根据本信息``激活预配置调度的预配置调度标识''和``去激活预配置调度的预配置调度标识''字段对应的高层配置，可能存在与本信息对应的一种或多种 TCI，例如本信息的单级 HARQ 反馈信息、本信息调度数据的单级 HARQ 反馈信息、本信息调度数据的两级 HARQ 反馈信息、本信息调度的 CQI 反馈信息等。若存在上述一种或多种 TCI，则上述一种或多种 TCI 都在本字段指示的基础载波上映射；若不存在上述任何一种 TCI，则本字段为无效字段。
  \item 6 比特：预留比特，当前版本全部为 0。
  \item 24 比特：根据本文件 6.5.1 节的处理方法，使用循环冗余校验生成多项式 $g_{\mathrm{CRC24B}}(D)$ 计算循环冗余校验序列，并增加 24 比特预配置调度激活去激活信息掩码。预配置调度激活去激活信息掩码 \texttt{newphy-IDforPreConfig} 由高层配置。
\end{itemize}

\subsubsection{6.2.4.6.3 休眠与唤醒指示信息}

休眠与唤醒指示信息的长度为 82 比特。从最低位比特到最高位比特，休眠与唤醒指示信息包括：
\begin{itemize}[nosep]
  \item 58 比特：每一比特指示特定 T 节点的休眠或者唤醒。当该比特为 1 时指示特定 T 节点在非连续接收的下一个持续时间内监听 G 链路控制信息，当比特为 0 时指示特定 T 节点在非连续接收的下一个持续时间内不监听 G 链路控制信息。58 个比特最大可指示 58 个特定 T 节点。T 节点根据高层信令的配置确定该 T 节点对应的比特位置。
  \item 24 比特：根据本文件 6.5.1 节的处理方法，使用循环冗余校验生成多项式 $g_{\mathrm{CRC24B}}(D)$ 计算循环冗余校验序列，并增加 24 比特休眠与唤醒指示信息掩码。休眠与唤醒指示信息掩码 \texttt{wakeupMASK} 由高层配置。
\end{itemize}

\subsubsection{6.2.4.6.4 快速载波切换指示信息}

快速载波切换指示信息中，有效比特 45 比特，预留比特 13 比特，CRC24 比特，共 82 比特。从最低位比特到最高位比特，快速载波切换指示信息包含如下信息：
\begin{itemize}[nosep]
  \item 4 比特：物理层控制信息格式指示。对于快速载波切换指示信息，该字段的取值为 2。
  \item 1 比特：切换工作载波的节点类型。0：G 节点，1：T 节点。
  \item 16 比特：切换前工作载波对应的信道的信道号。
  \item 16 比特：切换后工作载波对应的信道的信道号。
  \item 8 比特：切换时刻指示。以发送当前快速载波切换指示信息的 TTI 为起点，指示载波切换前，当前载波最后一个 TTI 的偏移量，单位为 TTI。取值为 0，表示当前发送快速载波切换指示信息的 TTI 即当前载波的最后一个 TTI。在非连续传输模式下，本信息的``切换前工作载波对应的信道的信道号''字段与``切换后工作载波对应的信道的信道号''字段的取值相同时，本字段指示的 TTI 为当前 COT 的最后一个 TTI。
  \item 13 比特：预留比特，该版本全部为 0。
  \item 24 比特：根据本文件 6.5.1 节的处理方法，使用循环冗余校验生成多项式 $g_{\mathrm{CRC24B}}(D)$ 计算循环冗余校验序列。当快速载波切换指示信息用于指示所有当前载波与本控制信道中``切换前载波对应的信道的信道号''字段一致的 T 节点切换载波时，或当工作载波的节点为 G 节点时，使用 24 比特 G 节点标识进行加扰；当快速载波切换指示信息用于指示特定 T 节点切换载波时，使用该 T 节点的 24 比特 T 节点标识进行加扰。
\end{itemize}

\subsubsection{6.2.4.6.5 初始接入载波指示信息}

初始接入载波指示信息中，有效比特 52 比特，预留比特 6 比特，CRC24 比特，共 82 比特。从最低位比特到最高位比特，初始接入载波指示信息包含如下信息：
\begin{itemize}[nosep]
  \item 4 比特：物理层控制信息格式指示。对于初始接入载波指示信息，该字段的取值为 3。
  \item 16 比特：工作载波包含 1 个基础载波的初始接入阶段 T 节点工作载波对应的信道的信道号。
  \item 16 比特：工作载波包含 4 个基础载波的初始接入阶段 T 节点工作载波对应的信道的信道号。
  \item 16 比特：工作载波包含本频段内所有基础载波的初始接入阶段 T 节点工作载波对应的信道的信道号。
  \item 6 比特：预留比特，该版本全部为 0。
  \item 24 比特：根据本文件 6.5.1 节的处理方法，使用循环冗余校验生成多项式 $g_{\mathrm{CRC24B}}(D)$ 计算循环冗余校验序列。使用 24 比特 G 节点标识进行加扰。
\end{itemize}

T 节点在初始接入过程中，根据本信息确定初始接入阶段 T 节点工作载波。T 节点初始接入过程中，在发送接入信息之前若未收到本信息，则在符合接入信息``指示初始接入阶段 T 节点带宽类型''字段描述的，且与初始接入信道对应的工作载波中，选择接入信息块所在的工作载波作为初始接入阶段 T 节点工作载波。初始接入信道的定义参见本标准 8.3 节。

若 G 节点配置了初始接入阶段 T 节点工作载波，则在每个包含接入信息块的基础载波上，每个包含接入信息块的同步信息块周期内，G 节点在其中至少 1 个 TTI 内发送本信息。G 节点发送本信息的资源不受同步信息中``指示本基础载波包含哪些 T 节点组的 GCI''字段的限制。

\subsubsection{6.2.4.6.6 G 链路位置测量信号或感知测量信号资源指示信息}

位置测量信号资源指示信息用于指示当前 TTI 内位置测量信号的时间资源，共 82 比特，支持四种不同格式。

\paragraph{第一格式（双边两信号）}

用于指示当前 TTI 内双边两信号的接收或发送，位置测量信号资源指示信息从最低位比特到最高位比特，包括 82 比特：
\begin{itemize}[nosep]
  \item 4 比特：物理层控制信息格式指示。1 指示位置测量资源配置。
  \item 3 比特：位置测量资源配置格式指示。该字段值为 0 时，使用如下字段的指示信息。
  \item 3 比特：预留，该版本全为 0。
  \item 1 比特：信号收发指示信息。0 指示使用第一资源发送测量信号，1 指示使用第一资源接收测量信号。
  \item 7 比特：第一资源起始测距符号指示信息。在被调度的 TTI 里，起始测距符号为 \#(该字段数值) 符号。对于小于或等于 1\,ms 的 TTI，其指示单位为 1 个符号；对于 2\,ms 的 TTI，其指示单位为 2 个符号；对于 4\,ms TTI，其指示单位为 4 个符号；对于 8\,ms TTI，其指示单位为 8 个符号。
  \item 4 比特：第一资源天线端口数量 $N_{G,\mathrm{port}}$，0 表示不进行资源配置。
  \item 3 比特：第一资源天线端口测距符号重复次数 $N_{G,\mathrm{port,rep}}$，假如该域值为 $v$，则一个端口发送的测距符号的个数为 $N_{G,\mathrm{port,rep}} = 2^v$。
  \item 5 比特：第一资源波束增强测距波束数量 $N_{G,\mathrm{bf}}$，假如该域值为 $v$，则 $N_{G,\mathrm{bf}} = 2^v$，0 表示不进行资源配置。
  \item 3 比特：第一资源波束增强测距符号重复次数 $N_{G,\mathrm{bf,rep}}$，假如该域值为 $v$，则一个波束发送的测距符号的个数为 $N_{G,\mathrm{bf,rep}} = 2^v$。
  \item 1 比特：第一资源波束增强测距切换间隔指示信息。0 指示不需要切换间隔符号，1 指示需要切换间隔符号。
  \item 1 比特：信号收发指示信息。0 指示使用第二资源发送测量信号，1 指示使用第二资源接收测量信号。
  \item 7 比特：第二资源起始测距符号指示信息。在被调度的 TTI 里，起始测距符号为 \#(该字段数值) 符号。对于小于或等于 1\,ms 的 TTI，其指示单位为 1 个符号；对于 2\,ms 的 TTI，其指示单位为 2 个符号；对于 4\,ms TTI，其指示单位为 4 个符号；对于 8\,ms TTI，其指示单位为 8 个符号。
  \item 4 比特：第二资源天线端口数量 $N_{T,\mathrm{port}}$，0 表示不进行资源配置。
  \item 3 比特：第二资源天线端口测距符号重复次数 $N_{T,\mathrm{port,rep}}$，假如该域值为 $v$，则一个端口发送的测距符号的个数为 $N_{T,\mathrm{port,rep}} = 2^v$。
  \item 5 比特：第二资源波束增强增强测距波束数量 $N_{T,\mathrm{bf}}$，假如该域值为 $v$，则 $N_{T,\mathrm{bf}} = 2^v$，0 表示不进行资源配置。
  \item 3 比特：第二资源波束增强测距符号重复次数 $N_{T,\mathrm{bf,rep}}$，假如该域值为 $v$，则一个波束发送的测距符号的个数为 $N_{T,\mathrm{bf,rep}} = 2^v$。
  \item 1 比特：第二资源波束增强测距切换间隔指示信息。0 指示不需要切换间隔符号，1 指示需要切换间隔符号。
  \item 24 比特：根据本标准 6.5.1 节的处理方法，使用循环冗余校验生成多项式 $g_{\mathrm{CRC24B}}(D)$ 计算循环冗余校验序列，并增加 24 比特位置测量信号资源指示信息掩码。24 比特位置测量信号资源指示信息掩码 \texttt{positioningResourceIndicationMASK} 由高层配置。
\end{itemize}

对于第一格式位置测量信号资源指示信息，当 $N_{G,\mathrm{port}}$ 和 $N_{G,\mathrm{bf}}$ 的值都为零时，第一资源指示信息无效，当 $N_{T,\mathrm{port}}$ 和 $N_{T,\mathrm{bf}}$ 的值都为 0 时，第二资源指示信息无效。当指示信息仅指示第一资源信息或第二资源信息时，指示信息用于指示单边单信号的接收或发送；当指示信息指示第一资源信息和第二资源信息时，指示信息用于指示双边两信号的接收和发送。根据第一资源信息和第二资源信息的时间先后顺序，以及第一资源信息和第二资源信息的发送和接收行为，可以配置 G 节点先发测量信号、T 节点后发测量信号，或者 T 节点先发测量信号、G 节点后发测量信号，或者一个 T 点先发测量信号、另外一个 T 节点后发测量信号。

\paragraph{第二格式（单端口双边三信号）}

用于指示当前 TTI 内单端口双边三信号的接收或发送，位置测量信号资源指示信息从最低位比特到最高位比特，包括 82 比特：
\begin{itemize}[nosep]
  \item 4 比特：物理层控制信息格式指示。1 指示位置测量资源配置。
  \item 3 比特：位置测量资源配置格式指示。该字段值为 1 时，使用如下字段的指示信息。
  \item 3 比特：预留，该版本全为 0。
  \item 1 比特：三信号资源接收发送指示，0 指示发送-接收-发送，1 指示接收-发送-接收。
  \item 7 比特：第一信号资源起始测距符号指示信息。在被调度的 TTI 里，起始测距符号为 \#(该字段数值) 符号。对于小于或等于 1\,ms 的 TTI，其指示单位为 1 个符号；对于 2\,ms 的 TTI，其指示单位为 2 个符号；对于 4\,ms TTI，其指示单位为 4 个符号；对于 8\,ms TTI，其指示单位为 8 个符号。
  \item 5 比特：第一信号资源单端口波束数量 $N_{m1,\mathrm{bf}}$，假如该域值为 $v$，则 $N_{m1,\mathrm{bf}} = 2^v$，0 表示不进行资源配置。
  \item 3 比特：第一信号资源单端口波束测距符号的重复次数 $N_{m1,\mathrm{bf,rep}}$，假如该域值为 $v$，则一个波束发送的测距符号的个数为 $N_{m1,\mathrm{bf,rep}} = 2^v$。
  \item 1 比特：第一信号切换间隔指示信息。0 指示不需要切换间隔符号，1 指示需要切换间隔符号。
  \item 7 比特：第二信号资源起始测距符号指示信息。在被调度的 TTI 里，起始测距符号为 \#(该字段数值) 符号。对于小于或等于 1\,ms 的 TTI，其指示单位为 1 个符号；对于 2\,ms 的 TTI，其指示单位为 2 个符号；对于 4\,ms TTI，其指示单位为 4 个符号；对于 8\,ms TTI，其指示单位为 8 个符号。
  \item 5 比特：第二信号资源单端口波束数量 $N_{m2,\mathrm{bf}}$，假如该域值为 $v$，则 $N_{m2,\mathrm{bf}} = 2^v$，0 表示不进行资源配置。
  \item 3 比特：第二信号资源单端口波束测距符号的重复次数 $N_{m2,\mathrm{bf,rep}}$，假如该域值为 $v$，则一个波束发送的测距符号的个数为 $N_{m2,\mathrm{bf,rep}} = 2^v$。
  \item 1 比特：第二信号切换间隔和第三信号切换间隔指示信息。0 指示第二信号和第三信号之前不需要切换间隔符号，1 指示第二信号和第三信号之前需要切换间隔符号。
  \item 7 比特：第三信号资源起始测距符号指示信息。在被调度的 TTI 里，起始测距符号为 \#(该字段数值) 符号。对于小于或等于 1\,ms 的 TTI，其指示单位为 1 个符号；对于 2\,ms 的 TTI，其指示单位为 2 个符号；对于 4\,ms TTI，其指示单位为 4 个符号；对于 8\,ms TTI，其指示单位为 8 个符号。
  \item 5 比特：第三信号资源单端口波束数量 $N_{m3,\mathrm{bf}}$，假如该域值为 $v$，则 $N_{m3,\mathrm{bf}} = 2^v$，0 表示不进行资源配置。
  \item 3 比特：第三信号资源单端口波束测距符号的重复次数 $N_{m3,\mathrm{bf,rep}}$，假如该域值为 $v$，则一个波束发送的测距符号的个数为 $N_{m3,\mathrm{bf,rep}} = 2^v$。
  \item 24 比特：根据本标准 6.5.1 节的处理方法，使用循环冗余校验生成多项式 $g_{\mathrm{CRC24B}}(D)$ 计算循环冗余校验序列，并增加 24 比特位置测量信号资源指示信息掩码。24 比特位置测量信号资源指示信息掩码 \texttt{positioningResourceIndicationMASK} 由高层配置。
\end{itemize}

\paragraph{第三格式（多端口双边三信号）}

用于指示当前 TTI 多端口双边三信号的接收或发送，位置测量信号资源指示信息从最低位比特到最高位比特，包括 82 比特：
\begin{itemize}[nosep]
  \item 4 比特：物理层控制信息格式指示。1 指示位置测量资源配置。
  \item 3 比特：位置测量资源配置格式指示。该字段值为 2 时，使用如下字段的指示信息。
  \item 2 比特：预留，该版本全为 0。
  \item 1 比特：三消息资源接收发送指示，0 指示发送-接收-发送，1 指示接收-发送-接收。
  \item 7 比特：第一信号资源起始测距符号指示信息。在被调度的 TTI 里，起始测距符号为 \#(该字段数值) 符号。对于小于或等于 1\,ms 的 TTI，其指示单位为 1 个符号；对于 2\,ms 的 TTI，其指示单位为 2 个符号；对于 4\,ms TTI，其指示单位为 4 个符号；对于 8\,ms TTI，其指示单位为 8 个符号。
  \item 7 比特：第三信号资源起始测距符号指示信息。在被调度的 TTI 里，起始测距符号为 \#(该字段数值) 符号。对于小于或等于 1\,ms 的 TTI，其指示单位为 1 个符号；对于 2\,ms 的 TTI，其指示单位为 2 个符号；对于 4\,ms TTI，其指示单位为 4 个符号；对于 8\,ms TTI，其指示单位为 8 个符号。
  \item 3 比特：第一/三信号资源天线端口数量 $N_{\mathrm{port}}$，0 表示不进行资源配置。
  \item 2 比特：第一/三信号资源天线端口测距重复次数 $N_{\mathrm{port,rep}}$，假如该域值为 $v$，则 $N_{\mathrm{port,rep}} = 2^{(v+1)}$。
  \item 5 比特：第一/三信号资源波束增强测距波束数量 $N_{\mathrm{bf}}$，假如该域值为 $v$，则 $N_{\mathrm{bf}} = 2^v$，0 表示不进行资源配置。
  \item 3 比特：第一/三信号资源波束增强测距波束重复次数 $N_{\mathrm{bf,rep}}$，假如该域值为 $v$，则一个波束发送的测距符号的个数为 $N_{\mathrm{bf,rep}} = 2^v$。
  \item 1 比特：第一/第二/第三信号波束增强测距切换间隔指示信息。0 指示不需要切换间隔符号，1 指示需要一个切换间隔符号。
  \item 7 比特：第二信号资源起始测距符号指示信息。在被调度的 TTI 里，起始测距符号为 \#(该字段数值) 符号。对于小于或等于 1\,ms 的 TTI，其指示单位为 1 个符号；对于 2\,ms 的 TTI，其指示单位为 2 个符号；对于 4\,ms TTI，其指示单位为 4 个符号；对于 8\,ms TTI，其指示单位为 8 个符号。
  \item 3 比特：第二信号资源天线端口数量 $N_{\mathrm{port}}$，0 表示不进行资源配置。
  \item 2 比特：第二信号资源天线端口测距重复次数 $N_{\mathrm{port,rep}}$，假如该域值为 $v$，则 $N_{\mathrm{port,rep}} = 2^{(v+1)}$。
  \item 5 比特：第二信号资源波束增强测距波束数量 $N_{\mathrm{bf}}$，假如该域值为 $v$，则 $N_{\mathrm{bf}} = 2^v$，0 表示不进行资源配置。
  \item 3 比特：第二信号资源波束增强测距波束重复次数 $N_{\mathrm{bf,rep}}$，假如该域值为 $v$，则一个波束发送的测距符号的个数为 $N_{\mathrm{bf,rep}} = 2^v$。
  \item 24 比特：根据本标准 6.5.1 节的处理方法，使用循环冗余校验生成多项式 $g_{\mathrm{CRC24B}}(D)$ 计算循环冗余校验序列，并增加 24 比特位置测量信号资源指示信息掩码。24 比特位置测量信号资源指示信息掩码 \texttt{positioningResourceIndicationMASK} 由高层配置。
\end{itemize}

根据指示信息，节点从起始测距符号开始，连续使用 $N_{\mathrm{port}} \times N_{\mathrm{port,rep}}$ 个符号发送基于端口的位置测量信号，每 $N_{\mathrm{port,rep}}$ 个符号一组，发送相同的位置测量信号。$N_{\mathrm{port}}$ 组位置测量信号使用 $N$ 点 DFT 矩阵进行扩频；根据指示信息，节点从起始测距符号开始，使用 $N_{\mathrm{bf}} \times N_{\mathrm{bf,rep}}$ 个符号发送基于波束的位置测量信号，当不需要切换间隔符号时，每个发送位置测量信号的符号连续，当需要切换间隔符号时，在使用连续 $N_{\mathrm{bf,rep}}$ 个符号发送 $N_{\mathrm{bf,rep}}$ 次波束后紧邻的一个符号不发送任何信息。

位置测量信号的带宽 \texttt{posFrequencyBandwidthIndication} 由高层配置。

\paragraph{第四格式（测量组激活/去激活）}

用于指示当前 TTI 内测量组的激活或去激活，位置测量信号资源指示信息从最低位比特到最高位比特，包括 82 比特：
\begin{itemize}[nosep]
  \item 4 比特：物理层控制信息格式指示。1 指示位置测量资源配置。
  \item 3 比特：位置测量资源配置格式指示。该字段值为 3 时，使用如下字段的指示信息。
  \item 6 比特：预留，该版本全为 0。
  \item 45 比特：每一个测量组激活去激活指示信息共 15 比特，其中前 8 比特是测量组 ID，指示有连接或无连接的标签区分是否是目标测量组的 GCI；第 9 比特是该资源配置的激活或者去激活指示，当该比特为 1 时是测量组测量激活指示，当比特为 0 时是测量组测量去激活指示。第 10--11 比特设置为 0 时，表示不调度测量时间差报告，设置为 1 时指示调度时间差报告，即请求测量组 ID 对应的测量组中的所有锚点和标签向 G 节点发送时间差测量报告消息，设置为 2 时，指示测量组 ID 对应的测量组的所有锚点依次向测量组中的标签发送时间差报告消息，且各锚点按照测量参数配置消息中规定的 MCS 和带宽发送时间差报告。第 12--15 比特预留，该版本全为 0。45 比特共指示 3 个测量组的位置测量资源配置激活或去激活信息。
  \item 24 比特：根据本标准 6.5.1 节的处理方法，使用循环冗余校验生成多项式 $g_{\mathrm{CRC24B}}(D)$ 计算循环冗余校验序列，并增加 24 比特测量组组播地址作为信息掩码，该测量组组播地址信息掩码由高层配置。
\end{itemize}

\subsubsection{6.2.4.6.7 T 链路位置测量信号资源指示信息}

T 节点发送 T 链路位置测量信号资源指示信息，指示测量组内的多个锚点按顺序发送测量时间差报告消息。上述 T 节点为测量组建立消息指示的第一个锚点。测量组建立消息指示的第二个锚点及后续锚点不发送 T 链路位置测量信号资源指示信息。

T 链路位置测量信号资源指示信息共 85 比特，从最低位比特到最高位比特，包括：
\begin{itemize}[nosep]
  \item 4 比特：用于每个位置锚点发送测量时间差报告消息的时间资源，以 TTI 为单位。
  \item 3 比特：解调参考信号端口数量指示信息。解调参考信号端口数量是（3 比特数值+1）。
  \item 16 比特：子载波组指示信息。16 比特中从最低位的比特开始到最高位的比特结束，与子载波组从小到大的顺序一一对应。16 比特中，取值为 1 的比特代表使用该比特对应的子载波组，取值为 0 的比特代表不使用该比特对应的子载波组。
  \item 10 比特：起始符号指示。在该 TTI 内，测量时间差报告的起始符号的索引。
  \item 10 比特：结束符号指示。在该 TTI 内，测量时间差报告的结束符号的索引。
  \item 2 比特：HARQ 进程指示。
  \item 1 比特：新包指示。当该值翻转时，指示传输新包，当该值不变时，指示重传旧包。
  \item 5 比特：调制编码方式指示信息。对于单 TB 传输，该字段指示该单 TB 的调制编码方式；对于两 TB 传输，该字段指示第一个 TB 的调制编码方式。
  \item 3 比特：调制符号重复指示。对于单流传输，该字段指示该单流传输的重复次数；对于两流传输，该字段指示第 2 流传输时符号的重复次数。
  \item 2 比特：0 指示重传版本 0，1 指示重传版本 1，2 指示重传版本 2，3 指示重传版本 3。
  \item 5 比特：在高层 XRC 信令配置的多个 ACK/NACK 资源池中指示资源。
  \item 24 比特：根据本文件 6.5.1 节的处理方法，使用循环冗余校验生成多项式 $g_{\mathrm{CRC24B}}(D)$ 计算循环冗余校验序列，并使用 24 比特测量组对应的组播地址物理层标识进行加扰。
\end{itemize}

\subsubsection{6.2.4.6.8 G 链路控制信息信道编码、调制和物理资源映射}

G 链路控制信息传输使用本文件 6.2.6 规定的信道编码参数，以 $C = 1$ 按照 6.2.6.1.2、6.2.6.1.3、6.2.6.1.4 进行物理层控制信息的极化编码，产生比特序列 $g_k$，其中 $k = 0, \ldots, G-1$，并根据本文件 6.5.6 节的方法使用 QPSK 调制将比特序列 $g_k$ 映射为调制符号序列 $d_i$。当使用 $L$ 个 CTU 进行传输时，调制符号序列 $d_i$ 在 CTU 的可用资源内按照先低频率子载波再高频率子载波，并按照符号的时间先后顺序进行映射。

\subsection{6.2.4.7 G 链路物理层控制信息块 G-PCIB}

G 链路物理层控制信息块用于 G 节点映射 G 链路控制信息，包括动态调度数据控制信息、预配置调度激活去激活信息、休眠与唤醒指示信息、快速载波切换指示信息、G 链路位置测量信号资源指示信息、T 链路位置测量信号资源指示信息。

在每一个包含同步信息块的基础载波上，每一个 TTI 中包含 NSymb-G-PCIB 个符号为一个 G-PCIB 的资源，NSymb-G-PCIB 在广播信息中指示，取值范围 $\{1, 2, 3, 4, 6, 8, 10, 12\}$。在包含同步信息块 SAB 和物理层广播信息的 TTI 中，G-PCIB 为从 \#8 符号开始的连续 NSymb-G-PCIB 个符号；在包含同步信息块 SAB 且不包含物理层广播信息的 TTI 中，若 NSymb-G-PCIB$>1$ 则 G-PCIB 为 \#5 符号和从 \#7 符号开始的连续 NSymb-G-PCIB$-1$ 个符号，若 NSymb-G-PCIB=1 则 G-PCIB 仅包含 \#5 符号；在其它 TTI 中，G-PCIB 为从 \#1 符号开始的连续 NSymb-G-PCIB 个符号。

在一个 G-PCIB 的符号中，包括 2 个 G 链路控制信息的最小传输资源单位（Control information transmission unit---CTU），其中 \#0--\#79 子载波定义为一个 CTU，\#81--\#160 子载波定义为一个 CTU。在 G 链路控制信息公共资源内，以及 T 节点被配置的 G 链路控制信息特定资源内，按照资源内 G 符号子载波从小到大的顺序，以及 G 符号时间的先后顺序，\#0 开始编号。以 1 个 CTU 为单位，G 链路控制信息聚合级别为 1、2、4 和 8。在 G 链路控制信息公共资源和 T 节点特定资源内，当一个控制信息使用聚合级别 $L$ 进行发送时，需要使用连续 $L$ 个 CTU，并且起始 CTU 索引 $i$ 满足 $\mod(i, L) = 0$。

按照同步信息的指示，在包含公共 GCI 的基础载波上，T 节点在每个 TTI 内的 G-PCIB 中盲检公共 GCI，T 节点在该 G-PCIB 中所有符合表 7 所示初始聚合级别的资源上盲检公共 GCI。

按照同步信息的指示，在包含至少一个 T 节点组 GCI 的基础载波上，属于其中一个 T 节点组的 T 节点在每个 TTI 内的 G-PCIB 中盲检 T 节点特定 GCI。在该 G-PCIB 中，可以为该 T 节点配置不大于 10 个盲检 T 节点特定 GCI 的资源。当未被配置盲检 T 节点特定 GCI 的资源时，该 T 节点在该 G-PCIB 中所有符合表 7 所示初始聚合级别的资源上盲检 T 节点特定 GCI。

T 节点应执行 G-PCIB 上所有配置的盲检，并按照所有盲检成功的 GCI 被调度和配置。

\begin{table}[h]
\centering
\caption{表 7\quad NSymb-G-PCIB 与初始聚合级别的对应关系}
\begin{tabular}{ccccc}
\toprule
NSymb-G-PCIB & 1 & 2、3 & 4、6 & 8、10、12 \\
\midrule
初始聚合级别 & 1、2 & 1、2、4 & 2、4、8 & 4、8 \\
\bottomrule
\end{tabular}
\end{table}

\subsection{6.2.4.8 直通链路动态调度数据控制信息}

直通链路动态调度数据控制信息有效比特 52 比特，预留比特 6 比特，CRC24 比特，共 82 比特，从最低位比特到最高位比特，包括：
\begin{itemize}[nosep]
  \item 19 比特：COT 起始无线帧索引号。
  \item 10 比特：起始符号指示。COT 内起始符号的索引。
  \item 10 比特：结束符号指示。COT 内结束符号的索引。
  \item 8 比特：数据包序号 $N_{\mathrm{TB}}$。
  \item 5 比特：调制编码方式指示信息。
  \item 6 比特：预留比特，当前版本全部为 0。
  \item 24 比特：根据本文件 6.5.1 节的处理方法，使用循环冗余校验生成多项式 $g_{\mathrm{CRC24B}}(D)$ 计算循环冗余校验序列，并使用 24 比特直通链路先发节点物理层标识进行加扰。
\end{itemize}

直通链路动态调度数据控制信息使用本文件 6.2.6 规定的信道编码参数，以 $C = 1$ 按照 6.2.6.1.1、6.2.6.1.2、6.2.6.1.3 和 6.2.6.1.4 进行第二类数据信息的极化编码，产生比特序列 $g_k$，其中 $k = 0, \ldots, G-1$，并根据本文件 6.5.6 节的方法使用 QPSK 调制将比特序列 $g_k$ 映射为调制符号序列 $d_i$。当使用 $L$ 个 CTU 进行传输时，调制符号序列 $d_i$ 在 CTU 的可用资源内按照先低频率子载波再高频率子载波，并按照符号的时间先后顺序进行映射。

\subsection{6.2.4.9 直通链路物理层控制信息块 P-PCIB}

直通链路物理层控制信道用于直通链路先发节点映射直通链路动态调度数据控制信息。按照直通链路同步信息的指示，由直通链路控制信息聚合级别 NCTU-PCI 和直通链路控制信息数量 NPCI 确定 P-PCIB 的符号数量 $N_{\mathrm{Symb-P-PCIB}} = N_{\mathrm{CTU-PCI}} \times \lceil N_{\mathrm{PCI}}/2 \rceil$。

在直通链路同步信息块的基础载波上，P-PCIB 的资源包含直通链路同步信息块的无线帧中 \#5 符号，若 NSymb-P-PCIB$>1$，P-PCIB 的资源还包含直通链路同步信息块的无线帧中 \#7 符号以及之后的连续 NSymb-P-PCIB$-2$ 个符号。

在 P-PCIB 上，直通链路先发节点把 NPCI 个聚合级别为 NCTU-PCI 的基于传输块 TB 的动态调度数据控制信息按照先频后时的顺序，依次映射在 P-PCIB 上。

\subsection{6.2.4.10 T 链路控制信息 TCI}

T 链路控制信息，包括单级 HARQ 反馈信息、两级 HARQ 反馈信息、CQI 反馈信息、调度请求信息。T 链路控制信息使用基于序列传输方式或基于数据传输方式传输，不同类型的 T 链路控制信息支持的传输方式不同。

\subsubsection{6.2.4.10.1 单级 HARQ 反馈信息}

按照高层 \texttt{harqTxMode} 配置，单级 HARQ 反馈信息采用基于序列传输方式传输或基于数据传输方式传输。当 T 节点在 TTI 接收 G 链路第二类数据信息时，T 节点按照高层配置，在该 TTI 或者之后的 TTI 内使用该数据对应的 ACK 反馈信息资源池反馈 T 链路 ACK 反馈信息。

单级 HARQ 反馈信息中，有效比特 $N_1 + N_2$ 比特，从最低位比特到最高位比特，包括：
\begin{itemize}[nosep]
  \item $N_1$ 比特：ACK/NACK 反馈。按照高层 \texttt{schedMode} 配置，当采用基于 TB 的反馈时，ACK/NACK 反馈的比特数 $N_1$ 等于 1，与 TB 对应的 ACK/NACK 信息对应；当采用基于 CBG 的反馈时，ACK/NACK 反馈的比特数 $N_1$ 由高层 \texttt{singleStageCBGNumber} 配置，每 1 个比特依次与 CBG\#0 到 CBG\#$(N_1-1)$ 的 ACK/NACK 反馈信息对应。对于本字段的每 1 比特，``1''表示 ACK，``0''表示 NACK。
  \item $N_2$ 比特：MCS 调整量。MCS 调整量的比特数 $N_2$ 由高层 \texttt{mcsOffsetFeedbackBits} 配置。
\end{itemize}

当采用基于 TB 反馈 ACK/NACK 时：
\begin{itemize}[nosep]
  \item 当反馈是 ACK 时，MCS 调整量比特对应的数值表示推荐的 MCS 值相对于当前 MCS 值提升 [（比特数值）$\times$ MCS 偏移量调整的步进值] 个等级，MCS 偏移量调整的步进值根据高层 \texttt{mcsOffsetAdjustGranularity} 配置；
  \item 当反馈是 NACK 时，MCS 调整量比特对应的数值表示推荐的 MCS 值相对于当前 MCS 值降低（比特数值）$\times$ MCS 偏移量调整的步进值] 个等级，MCS 偏移量调整的步进值根据高层 \texttt{mcsOffsetAdjustGranularity} 配置。
\end{itemize}

当采用基于 CBG 反馈 ACK/NACK 时：
\begin{itemize}[nosep]
  \item 当所有 CBG 都反馈 ACK 时，MCS 调整量比特对应的数值表示推荐的 MCS 值相对于当前 MCS 值提升 [（比特数值）$\times$ MCS 偏移量调整的步进值] 个等级，MCS 偏移量调整的步进值根据高层 \texttt{mcsOffsetAdjustGranularity} 配置；
  \item 当至少有一个 CBG 反馈 NACK 时，MCS 调整量比特对应的数值表示推荐的 MCS 值相对于当前 MCS 值降低（比特数值）$\times$ MCS 偏移量调整的步进值] 个等级，MCS 偏移量调整的步进值根据高层 \texttt{mcsOffsetAdjustGranularity} 配置。
\end{itemize}

当采用基于序列传输方式传输单级 HARQ 反馈信息时，单级 HARQ 反馈信息仅包含有效比特。

当采用基于数据传输方式传输单级 HARQ 反馈信息时，单级 HARQ 反馈信息中，有效比特 $N_1 + N_2$ 比特，CRC24 比特，共 $N_1 + N_2 + 24$ 比特，从最低位比特到最高位比特，包括：
\begin{itemize}[nosep]
  \item $N_1 + N_2$ 比特：有效比特，参加本节前文单级 HARQ 反馈信息中有效比特的定义。
  \item 24 比特：根据本文件 6.5.1 节的处理方法，使用循环冗余校验生成多项式 $g_{\mathrm{CRC24B}}(D)$ 计算循环冗余校验序列。
\end{itemize}

\subsubsection{6.2.4.10.2 两级 HARQ 反馈信息}

两级 HARQ 反馈信息支持基于数据传输方式传输。

第一级反馈信息中，有效比特 $N_{S1}$ 比特（$N_{S1}$ 由高层 \texttt{firstStageBits} 指示），CRC24 比特，共 $N_{S1} + 24$ 比特，从最低位比特到最高位比特，包括：
\begin{itemize}[nosep]
  \item $N_{S1} - Z$ 比特：第一级基于 CBG 的 ACK/NACK 反馈。第一级反馈的比特数 $N_{S1}$ 即 CBG 个数，由高层 \texttt{firstStageBits} 配置。对于其中每 1 个比特，``1''表示 ACK，``0''表示 NACK。第一级反馈中，每个 CBG 对应的 CB 根据初传时 TB 到 CBG 的映射确定。在一次传输中，一个 TB 包含的 CB 个数为 $C$，则该 TB 实际包含的 CBG 的数为 $M = \min(C, N_{S1})$。$M_1 = \mod(C, M)$，如果 $M_1 > 0$，$K_1 = \lceil C/M \rceil$，$K_2 = \lfloor C/M \rfloor$，在 $M$ 个 CBG 中，\#0 到 \#$(M_1-1)$ 的 CBG 包含 $K_1$ 个 CB，\#$M_1$ 到 \#$(M-1)$ 的 CBG 包含 $K_2$ 个 CB；如果 $M_1 = 0$，每个 CBG 包含 $(C/M)$ 个 CB。
  \item $Z$ 比特：MCS 调整量。MCS 调整量比特数 $Z$ 由高层 \texttt{mcsOffsetFeedbackBits} 配置。当所有 CBG 都反馈 ACK 时，数值表示推荐的 MCS 值相对于当前 MCS 值提升 [（比特数值）$\times$ MCS 偏移量调整的步进值] 个等级；当至少有一个 CBG 反馈 NACK 时，数值表示推荐的 MCS 值相对于当前 MCS 值降低（比特数值）$\times$ MCS 偏移量调整的步进值] 个等级，MCS 偏移量调整的步进值根据高层 \texttt{mcsOffsetAdjustGranularity} 配置。
  \item 24 比特：根据本文件 6.5.1 节的处理方法，使用循环冗余校验生成多项式 $g_{\mathrm{CRC24B}}(D)$ 计算循环冗余校验序列。
\end{itemize}

第二级反馈信息中，有效比特 $N_{S2}$ 比特（$N_{S2}$ 由高层 \texttt{secondStageBits} 指示），CRC24 比特，共 $N_{S2} + 24$ 比特，从最低位比特到最高位比特，包括：
\begin{itemize}[nosep]
  \item $N_{S2}$ 比特：第二级基于子 CBG 的 ACK/NACK 反馈。第二级反馈中，每个子 CBG 对应的 CB 根据第一级反馈指示的错误的 CBG 包含的 CB 到子 CBG 的映射确定。根据第一级反馈，一个 TB 中错误的 CBG 包含的 CB 个数为 $C$，则该错误的 CBG 实际包含的子 CBG 的个数为 $M = \min(C, N_{S2})$。$M_1 = \mod(C, M)$，如果 $M_1 > 0$，$K_1 = \lceil C/M \rceil$，$K_2 = \lfloor C/M \rfloor$，在 $M$ 个子 CBG 中，\#0 到 \#$(M_1-1)$ 的子 CBG 包含 $K_1$ 个 CB，\#$M_1$ 到 \#$(M-1)$ 的子 CBG 包含 $K_2$ 个 CB；如果 $M_1 = 0$，每个子 CBG 包含 $(C/M)$ 个 CB。
  \item 24 比特：根据本文件 6.5.1 节的处理方法，使用循环冗余校验生成多项式 $g_{\mathrm{CRC24B}}(D)$ 计算循环冗余校验序列。
\end{itemize}

\subsubsection{6.2.4.10.3 CQI 反馈信息}

CQI 反馈信息支持基于数据传输方式传输。

CQI 反馈信息按照高层 \texttt{CQI-FeedbackMode} 配置，采用基于基础载波的反馈方式或基于工作子载波组的反馈方式。

当配置为基于基础载波的反馈方式时，CSI 反馈信息中，有效比特 $4 \times N_{\mathrm{CH}}$ 比特，CRC16 比特，共 $4 \times N_{\mathrm{CH}} + 16$ 比特，从最低位比特到最高位比特，包括：
\begin{itemize}[nosep]
  \item $4 \times N_{\mathrm{CH}}$ 比特：信道质量指示信息。按照频率从低到高的顺序，每 4 比特依次指示 G 节点工作载波中一个基础载波的参考调制编码方式，参考调制编码方式为当前使用的调制编码方式表格中调制编码索引为 $2M+1$ 的调制编码方式，其中 $M$ 为该 4 比特对应的数值。
  \item 16 比特：根据本文件 6.5.1 节的处理方法，使用循环冗余校验生成多项式 $g_{\mathrm{CRC16}}(D)$ 计算循环冗余校验序列。
\end{itemize}

当配置为基于工作子载波组的反馈方式时，CSI 反馈信息中，有效比特 $64 \times N_{\mathrm{CH}}/L_{\mathrm{CH}}$ 比特，CRC16 比特，共 $64 \times N_{\mathrm{CH}}/L_{\mathrm{CH}} + 16$ 比特，从最低位比特到最高位比特，包括：
\begin{itemize}[nosep]
  \item $64 \times N_{\mathrm{CH}}/L_{\mathrm{CH}}$ 比特：信道质量指示信息。按照频率从低到高的顺序，每 4 比特依次指示 G 节点工作载波中一个 G 节点工作子载波组的参考调制编码方式，参考调制编码方式为当前使用的调制编码方式表格中调制编码索引为 $2M+1$ 的调制编码方式，其中 $M$ 为该 4 比特对应的数值。
  \item 16 比特：根据本文件 6.5.1 节的处理方法，使用循环冗余校验生成多项式 $g_{\mathrm{CRC16}}(D)$ 计算循环冗余校验序列。
\end{itemize}

本节中 $N_{\mathrm{CH}}$ 和 $L_{\mathrm{CH}}$ 的定义参见 6.2.2.4 节。

按照高层 \texttt{ttiPeriod} 和 \texttt{ttiOffset} 的配置，CQI 反馈信息可以配置为每次激活后单次传输或每次激活后在 COT 内周期性多次传输。

\subsubsection{6.2.4.10.4 调度请求信息}

调度请求信息支持基于数据传输方式传输。

调度请求信息中，有效比特 24 比特，预留比特 5 比特，CRC16 比特，共 45 比特，从最低位比特到最高位比特，包括：
\begin{itemize}[nosep]
  \item 24 比特：T 节点标识。
  \item 5 比特：预留比特，当前版本全部为 0。
  \item 16 比特：根据本文件 6.5.1 节的处理方法，使用循环冗余校验生成多项式 $g_{\mathrm{CRC16}}(D)$ 计算循环冗余校验序列。
\end{itemize}

按照高层 \texttt{ttiPeriod} 和 \texttt{ttiOffset} 的配置，在非连续传输模式下，可以发送调度请求信息的资源可以配置为在每个 COT 内、在每个包含 SAB 的基础载波上以不大于同步信息块传输周期的周期出现；在连续传输模式下，可以发送调度请求信息的资源可以配置为在任意一个或多个基础载波上周期出现。

若配置了可以发送调度请求信息的资源，T 节点发送调度请求的功能既可以通过发送调度请求信息实现，又可以通过发送接入信息实现；若不配置可以发送调度请求信息的资源，T 节点发送调度请求的功能通过发送接入信息实现，参见 6.2.4.4 节。

\subsubsection{6.2.4.10.5 T 链路控制信息信道编码和调制}

基于数据传输方式传输的 T 链路控制信息使用本文件 6.2.6 规定的信道编码参数，以 $C = 1$ 按照 6.2.6.1.2、6.2.6.1.3、6.2.6.1.4 进行物理层控制信息的极化编码，产生比特序列 $g_k$，其中 $k = 0, \ldots, G-1$。

除两级 HARQ 反馈信息中的第二级 HARQ 反馈信息外，其它基于数据传输方式传输的 T 链路控制信息根据本文件 6.5.6 节的方法使用 QPSK 调制将比特序列 $g_k$ 映射为调制符号序列 $d_i$。

当对应的 G 链路数据传输的调制方式为 QPSK、16QAM 时，两级 HARQ 反馈信息中的第二级 HARQ 反馈信息的调制方式采用 QPSK；当对应的 G 链路数据传输的调制方式高于 16QAM 时，两级 HARQ 反馈信息中的第二级 HARQ 反馈信息的调制方式采用 16QAM。根据本文件 6.5.6 节的方法将两级 HARQ 反馈信息中的第二级 HARQ 反馈信息的比特序列 $g_k$ 映射为调制符号序列 $d_i$。

\subsubsection{6.2.4.10.6 T 链路控制信息物理资源映射}

对于除调度请求信息外的其它 TCI，TCI 的资源由高层配置，动态调度数据控制信息或预配置调度激活去激活信息触发。对于除调度请求信息外的其它 TCI，映射 TCI 的基础载波，由对应的动态调度数据控制信息或预配置调度激活去激活信息中``TCI 所在的工作信道内的基础载波序号''字段指示。

调度请求信息由高层配置可用资源池，T 节点在有需求时自行选择资源池中的资源发送调度请求信息。T 节点在有需求时以概率 $P$ 在选择的资源中发送调度请求信息；T 节点可以多次重新选择资源并尝试以概率 $P$ 在重新选择的资源中发送调度请求信息，直到成功发送调度请求信息。$P$ 的初始值为 1；每当 T 节点成功发送调度请求信息后，在预设的时间窗内（当前版本预设的时间窗长由 T 节点自己确定）未收到对应的响应信息，$P$ 更新为当前值的一半，并可以再次以更新的 $P$；每当 T 节点收到对应的响应信息，$P$ 更新为 1。

TCI 在一个基础载波上的可用梳齿集合由高层配置。传输 T-DMRS-C 的频域资源（包括基础载波和梳齿）与传输对应的 TCI 的频域资源相同。

定义 TCI 符号索引：基准 TTI 中的第一个符号的符号索引为 0，之后每个符号的符号索引加 1。除调度请求信息外，其它 TCI 以发送对应的动态调度数据控制信息或预配置调度激活去激活信息的 TTI 为基准 TTI。调度请求信息以配置的可以发送调度请求信息的 TTI 为基准 TTI。

TCI 的发送符号中第一个符号对应的 TCI 符号索引为 $K_1$，从该符号开始的连续 $K_2$ 个符号为 TCI 的发送符号，$K_1$ 和 $K_2$ 分别由高层 \texttt{DedicatedTimeResource} 中的 \texttt{symbolStarting} 和 \texttt{symbolNumber} 指示。T 节点在 TCI 的发送符号中第一个符号上映射 T-DMRS-C，并在 TCI 的发送符号中其他符号上映射 TCI 信息。

对于基于数据传输方式传输的 TCI，T 节点在 TCI 的发送符号中的其它符号上按照符号的时间顺序，在每个符号中可用梳齿集合上，按照先低频率子载波再高频率子载波的顺序映射 TCI 的调制符号序列 $d_i$。两级 HARQ 反馈信息作为基于数据传输方式传输的 TCI 的一个特例，T 节点在 TCI 的发送符号中，从第二个符号开始的连续 $K_3$ 个符号上按照符号的时间顺序，在每个符号中可用梳齿集合上，按照先低频率子载波再高频率子载波的顺序映射第一级反馈信息的调制符号序列；并在之后的符号上按照符号的时间顺序，在每个符号中可用梳齿集合上，按照先低频率子载波再高频率子载波的顺序映射第二级反馈信息的调制符号序列，其中 $K_3$ 由高层 \texttt{firstStageSymbolNumber} 配置。

对于基于序列传输方式传输的 TCI，T 节点在 TCI 的发送符号中的其它符号上按照符号的时间顺序，在每个符号中按照可用梳齿集合中梳齿偏移序号从小到大的顺序对应信息比特从低位到高位的顺序，每个梳齿映射 TCI 中 1 比特信息对应的序列。对于基于序列传输方式传输的 TCI，应满足 $N \leq C \times (K_2 - 1)$，其中 $N$ 为 TCI 承载的比特数，$C$ 为可用梳齿集合中梳齿的数量。若一个超帧中 \#$n$ 无线帧中资源元素 $(k, l)$ 属于一个映射 TCI 中 1 比特信息的梳齿，在该梳齿上映射``1''对应的序列时，在该资源元素上映射的复数值 $a_{k,l} = r_{n,l}(k)$；在该梳齿上映射``0''对应的序列时，在该资源元素上映射的复数值 $a_{k,l} = -r_{n,l}(k)$。其中 $r_{n,l}(k)$ 为伪随机 QPSK 序列，生成方式见 6.5.3 节。

传输 T-DMRS-C 和对应的 TCI 的功率应按照在一个基础载波上的可用梳齿集合中的子载波占有效子载波的比例做反比例的功率推升（power boosting）。

%----------------------------------------------------------------------
\section{17 类控制信息对照表（6.2.1 索引 $\rightarrow$ 6.2.4 条款）}
%----------------------------------------------------------------------

\begin{longtable}{p{4.5cm}p{2.2cm}p{2.2cm}p{4.5cm}}
\caption{物理层 17 类控制信息与 6.2.4 子节对应} \\
\toprule
\textbf{控制信息（6.2.1）} & \textbf{6.2.4 子节} & \textbf{比特长度} & \textbf{主要承载 PIB/散布} \\
\midrule
\endfirsthead
\multicolumn{4}{c}{\tablename\ \thetable\ （续）} \\
\toprule
\textbf{控制信息} & \textbf{子节} & \textbf{长度} & \textbf{PIB} \\
\midrule
\endhead
同步信息 & 6.2.4.1 & 80 & SAB（\#3、\#4） \\
广播信息 & 6.2.4.3 & 72 & 散布（\#5、\#7） \\
接入信息 & 6.2.4.4 & 80 & RAB（\#3、\#4） \\
动态调度数据控制信息 & 6.2.4.6.1 & 82/88/91 & G-PCIB \\
预配置调度激活去激活信息 & 6.2.4.6.2 & 82 & G-PCIB \\
休眠与唤醒指示信息 & 6.2.4.6.3 & 82 & G-PCIB \\
快速载波切换指示信息 & 6.2.4.6.4 & 82 & G-PCIB \\
初始接入载波指示信息 & 6.2.4.6.5 & 82 & G-PCIB \\
G 链路位置/感知测量信号资源指示 & 6.2.4.6.6 & 82 & G-PCIB \\
T 链路位置测量信号资源指示 & 6.2.4.6.7 & 85 & G-PCIB \\
直通链路动态调度数据控制信息 & 6.2.4.8 & 82 & P-PCIB \\
单级 HARQ 反馈信息 & 6.2.4.10.1 & $N_1{+}N_2$（$+24$） & TCI 散布 \\
两级 HARQ 反馈信息 & 6.2.4.10.2 & $N_{S1}{+}24$ / $N_{S2}{+}24$ & TCI 散布 \\
CQI 反馈信息 & 6.2.4.10.3 & 变长 + CRC16 & TCI 散布 \\
调度请求信息 & 6.2.4.10.4 & 45 & TCI 散布 \\
\bottomrule
\end{longtable}

注：SAB（6.2.4.2）、RAB（6.2.4.5）、G-PCIB（6.2.4.7）、P-PCIB（6.2.4.9）为 PIB 容器，本身不单独计入 17 类控制信息清单。

%----------------------------------------------------------------------
\section{PIB 映射关系}
%----------------------------------------------------------------------

\begin{table}[h]
\centering
\footnotesize
\begin{tabular}{p{2.2cm}p{1.8cm}p{5.5cm}p{4cm}}
\toprule
\textbf{PIB} & \textbf{符号} & \textbf{块内控制/信号内容} & \textbf{关联 6.2.4 子节} \\
\midrule
SAB（通信域） & \#0--\#4 & STS、FTS、通信域同步信息 & 6.2.4.1、6.2.4.2 \\
SAB（直通） & \#0--\#4 & STS、FTS、直通同步信息 & 6.2.4.1、6.2.4.2 \\
---（广播） & \#5、\#7 & 广播信息（散布，非整包 PIB） & 6.2.4.3 \\
RAB & \#0--\#4（+可选 \#5、\#6） & STS、FTS、接入信息、随路短数据 & 6.2.4.4、6.2.4.5 \\
G-PCIB & 见 NSymb 规则 & GCI（DSCI/预配置/DRX/载波/定位等） & 6.2.4.6、6.2.4.7 \\
P-PCIB & \#5 起 & 直通 DSCI $\times$ NPCI & 6.2.4.8、6.2.4.9 \\
TCI & 高层配置符号 & HARQ/CQI/SR + T-DMRS-C & 6.2.4.10 \\
\bottomrule
\end{tabular}
\end{table}

\begin{table}[h]
\centering
\small
\caption{G-PCIB 符号位置与 TTI 类型}
\begin{tabular}{lll}
\toprule
\textbf{TTI 类型} & \textbf{G-PCIB 符号位置} & \textbf{备注} \\
\midrule
含 SAB + 广播 & 从 \#8 起连续 NSymb-G-PCIB 个 & 广播在 \#5、\#7 \\
含 SAB、无广播，NSymb$>1$ & \#5 + 从 \#7 起 NSymb$-1$ 个 & --- \\
含 SAB、无广播，NSymb=1 & 仅 \#5 & --- \\
其它 TTI & 从 \#1 起连续 NSymb-G-PCIB 个 & 常规业务 TTI \\
\bottomrule
\end{tabular}
\end{table}

\begin{table}[h]
\centering
\small
\caption{SAB/RAB 内 FTS $u$ 取值速查（与 6.2.3/6.2.4 衔接）}
\begin{tabular}{llll}
\toprule
\textbf{场景} & \textbf{STS $u$} & \textbf{FTS $u$} & \textbf{控制信息} \\
\midrule
通信域 SAB & G-ID 低 4 位 $+1$ & 1 或 160（符号类型） & 通信域同步信息 \\
直通 SAB & 先发节点 ID 低 4 位 $+1$ & 159 & 直通同步信息 \\
RAB 接入 & G-ID 低 4 位 $+1$ & 2 & 接入信息 \\
\bottomrule
\end{tabular}
\end{table}

%----------------------------------------------------------------------
\section{约束与实现要点}
%----------------------------------------------------------------------

\begin{enumerate}[nosep]
  \item \textbf{比特顺序}：所有控制信息字段从 LSB 到 MSB 排列，CRC 位于最高位侧；空口映射遵循 6.1 约定。
  \item \textbf{信息格式指示}：通信域同步信息=0，直通同步信息=1，接入信息=2；GCI 各类型另有格式指示（DSCI TB=0，快速载波切换=2，初始接入载波=3，位置测量=1 等）。
  \item \textbf{G-PCIB 盲检}：起始 CTU 索引 $i$ 须满足 $\mod(i,L)=0$；聚合级别 $L \in \{1,2,4,8\}$；表 7 限定 NSymb-G-PCIB 与初始聚合级别组合。
  \item \textbf{公共/特定 GCI}：同步信息 8\,bit GCI 位图指示本载波公共 GCI 与各 T 节点组 GCI；特定 GCI CRC 用 T 节点 Phy-ID 加扰。
  \item \textbf{DSCI 多格式}：TB（82\,bit）、CBG（88\,bit）、两级（91\,bit$\times$2）互斥，由高层 \texttt{schedMode}/\texttt{harqFeedBackMode} 决定；XRC 重配置须更新 Phy-ID。
  \item \textbf{预配置调度}：非连续模式 COT 结束自动去激活；激活/去激活 ID 字段互斥指示未配置 ID。
  \item \textbf{接入退避}：RAB 与 SR 均使用概率 $P$ 半衰退避（初始 $P=1$，无响应则减半，有响应则复位）。
  \item \textbf{TCI 传输方式}：单级 HARQ 支持序列/数据两种方式；两级 HARQ/CQI/SR 仅数据方式；序列方式须满足 $N \leq C(K_2-1)$。
  \item \textbf{TCI 功率}：T-DMRS-C 与 TCI 同梳齿，按占用有效子载波比例反比 power boosting。
  \item \textbf{直通 P-PCIB}：$N_{\mathrm{Symb-P-PCIB}} = N_{\mathrm{CTU-PCI}} \times \lceil N_{\mathrm{PCI}}/2 \rceil$；NPCI 个 DSCI 先频后时映射。
  \item \textbf{广播 NSymb-G-PCIB}：广播信息中 3\,bit 字段同时决定 G-PCIB 符号数与位置，实现须与 6.2.4.7 规则一致。
  \item \textbf{位置测量格式}：6.2.4.6.6 格式指示 0/1/2/3 对应四种子格式，掩码/CRC 加扰方式因格式而异（格式 4 用测量组组播地址）。
\end{enumerate}

%----------------------------------------------------------------------
\section{与标准其他章节的衔接}
%----------------------------------------------------------------------

\begin{itemize}[nosep]
  \item \textbf{6.2.1}：17 类控制信息分类索引与 PIB 定义，6.2.4 为各类字段的规范来源；
  \item \textbf{6.2.2}：RE 网格、TTI 长度、符号类型；广播/同步/接入的符号索引（\#3、\#5 等）均引用 6.2.2 帧结构；
  \item \textbf{6.2.3}：SAB/RAB 内 FTS/STS 序列；G-DMRS-C/P-DMRS-C/T-DMRS-C 与控制 RE 绑定；
  \item \textbf{6.2.5}：DSCI/预配置 GCI 调度第二类数据；G 链路调度资源须扣除 G-PCIB/SAB/广播等开销 RE；
  \item \textbf{6.2.6/6.5.1/6.5.6}：控制信息极化编码（$C=1$）、CRC24B/CRC16、QPSK 调制；
  \item \textbf{6.6/6.7}：G/T 位置测量信号资源指示触发 PMS/测量报告流程；
  \item \textbf{7.2.2.1/7.2.2.2}：随机接入（RAB/接入信息）与资源请求（SR/接入信息）媒体接入层流程；
  \item \textbf{7.5.3}：ACK 资源池、SR-Config、CQI-Conf、ControlResource 等高层参数与 6.2.4.10/6.2.4.6 字段对应；
  \item \textbf{第 8.3 节}：初始接入信道定义，与 6.2.4.6.5 初始接入载波指示配合。
\end{itemize}

\vfill
\begin{center}
\small\textcolor{gray}{精确参数与字段定义以 T/XS 10001-2025 正式标准为准；本文档仅供 SLB 120\,kHz 物理层实现与联调参考。}
\end{center}

\end{document}
